\subsection{Introduction to Polynomial Inequalities} 
\begin{itemize}
\item To solve a polynomial equation, we start with a factored polynomial equal to \makeblanksmall.
\item We do this because if a \makeblank[1in] is 0, then one of the \makeblank[1in] is also 0.
\item There are also facts about how products to compare to zero.
\item If a number is \makeblank 0, it is \makeblank[1.5in].
\item If a number is \makeblank 0, it is \makeblank[1.5in].
\item We can predict the \makeblank[1in] of a product from the \makeblank[1in] of the factors.
\end{itemize}
\begin{Example}
Solve the polynomial inequality
$$\left(x-3\right)\left(x+1\right)\geq 0$$
Note that the equation $\left(x-3\right)\left(x+1\right)=0$ has two solutions: $x=\makeblanksmall$ and $x=\makeblanksmall$. These divide the number line into parts:
\begin{lpic}{}{17}
\foreach \y/\lbl in {-1/{x<-1},-4/{x=-1},-7/{-1<x<3},-10/{x=3},-13/{x>3}}
	\regtextleft {1}{\y}{$\lbl$};
\regtextleft{1}{-16}{The solution set of the inequality is:};
\regtextleft{1}{-17}{or in interval notation:};
\end{lpic}
\end{Example}